This section fixes baseline notation and separates classical definitions from manuscript-specific provisional objects.

\subsection*{Classical definitions (standard)}

\begin{definition}[Completed zeta function]
\[
\xi(s) := \tfrac12 s(s-1)\pi^{-s/2}\Gamma(s/2)\zeta(s).
\]
We use the standard functional-equation symmetry $\xi(s)=\xi(1-s)$.
\end{definition}

\begin{definition}[Critical strip and critical line]
The critical strip is
\[
\{s\in\C:0<\RePart(s)<1\},
\]
and the critical line is
\[
\{s\in\C:\RePart(s)=\tfrac12\}.
\]
\end{definition}

\begin{definition}[Reflection involution]
Let $\sigma:\C\to\C$ be $\sigma(s)=1-s$. This map is an involution, and its fixed-point set is the critical line.
\end{definition}

\subsection*{Manuscript-specific provisional definitions}

\begin{definition}[Proposed coherence object $H$ \textnormal{(provisional definition; placeholder object)}]
\placeholder{A full operator-level construction is still required: ambient space, domain, action, and well-posedness properties.}
The symbol $H$ denotes a proposed spectral object used by the manuscript as a formalization target for the coherence program.
At the current stage, $H$ is reserved notation, not yet a fully constructed theorem-level object.
\end{definition}

\begin{definition}[Admissible point and admissible class \textnormal{(provisional definition; formalization target)}]
\placeholder{A full well-posedness theorem for the regularized trace below is still required.}
Fix the proposed coherence object $H$.
An \emph{admissible point} is a complex number $s\in\C$ such that both regularized resolvents
\[
(H-sI)^{-1}_{\mathrm{reg}},
\qquad
(H-(1-s)I)^{-1}_{\mathrm{reg}}
\]
exist in the chosen operator model and have finite regularized traces.
The \emph{admissible class} is
\[
\mathcal{A}:=\left\{s\in\C:
\Tr_{\mathrm{reg}}(H-sI)^{-1}\ \text{and}\ \Tr_{\mathrm{reg}}(H-(1-s)I)^{-1}\ \text{are both defined and finite}
\right\}.
\]
\end{definition}

\begin{definition}[Canonical balance functional \texorpdfstring{$\Phi_H$}{Phi_H} \textnormal{(single main-text candidate; conditional/provisional)}]
\placeholder{The regularization scheme and its theorem-level properties remain open proof obligations.}
For $s\in\mathcal{A}$, define
\[
\Phi_H(s)
:=
\Tr_{\mathrm{reg}}(H-sI)^{-1}
-
\Tr_{\mathrm{reg}}(H-(1-s)I)^{-1}.
\]
Thus $\Phi_H:\mathcal{A}\to\C$ is the \emph{regularized resolvent reflection-defect functional} of $H$.
The manuscript uses the single balance predicate
\[
\text{Bal}_H(s) \iff \Phi_H(s)=0 \qquad (s\in\mathcal{A}).
\]
Interpretation: $\Phi_H(s)=0$ encodes exact reflection balance of the regularized resolvent response at the pair $(s,1-s)$.
This is the canonical main-text candidate for the balance object, and it remains a theorem target in the current manuscript version.
\end{definition}

\begin{remark}[Why this is the canonical candidate in the main flow]
The functional above is used as the single main-text candidate because it directly measures the $s\leftrightarrow 1-s$ defect at operator level, which is the symmetry direction needed by the proof spine.
No alternative balance functional is used in the core manuscript flow.
\end{remark}

\begin{remark}[Status boundary for theorem-level use]
The definitions of $H$, admissibility, and $\Phi_H$ above are provisional.
Open bridges include: construction of the operator model, choice and invariance of regularization, existence/finite-value proofs on $\mathcal{A}$, and compatibility of this $\Phi_H$ with the zero-correspondence target.
No theorem-level claim may rely on these terms without first supplying those formal components.
\end{remark}

\begin{remark}[Traceability requirement]
When the provisional objects are upgraded, each promoted definition should include: (i) explicit hypotheses, (ii) well-posedness statement, and (iii) dependency links to the lemmas/theorems that justify subsequent use.
\end{remark}

\subsection*{Stage-6 provisional transport-defect interface}

\begin{definition}[No-drift defect functional \texorpdfstring{$\Delta_{\mathrm{nd}}$}{Delta\_nd} \textnormal{(provisional; audit/planning object)}]
Work under one synchronized package $\mathcal{H}_{\mathrm{S6}}^{\min}$, and fix:
\begin{enumerate}
  \item a Stage-5 hypothesis tuple $\Theta_{\mathrm{S5}}=(\theta_1,\dots,\theta_m)$;
  \item a transfer class $\mathfrak{C}_{\mathrm{tr}}$ and contradiction-relevant subclass
  $\mathfrak{C}_{\mathrm{tr}}^{\mathrm{ctr}}\subseteq\mathfrak{C}_{\mathrm{tr}}$;
  \item for each $C\in\mathfrak{C}_{\mathrm{tr}}$, a transfer-position set $U(C)\subseteq\mathcal{A}_{\mathrm{strip}}$;
  \item for each clause $\theta_j$ and each $C$, a clause-defect gauge
  $\delta_j^{(C)}:U(C)\to[0,\infty]$.
\end{enumerate}
Define the pointwise tuple-defect gauge
\[
d_{\Theta}^{(C)}(u):=\max_{1\le j\le m}\delta_j^{(C)}(u),
\]
and define the no-drift defect of $C$ by
\[
\Delta_{\mathrm{nd}}(C;\Theta_{\mathrm{S5}},\delta):=
\sup_{u\in U(C)}d_{\Theta}^{(C)}(u)\in[0,\infty].
\]
Thus $\Delta_{\mathrm{nd}}$ is a \emph{transport-instability measure}: a supremum defect over positions $u$ and clauses $\theta_j$ along a transfer object $C$.
Intended semantics (not proved): $\Delta_{\mathrm{nd}}(C)=0$ means exact clause preservation along $C$; $0<\Delta_{\mathrm{nd}}(C)\ll1$ means quantitatively small drift; finite uniform bounds on a transfer subclass mean bounded instability; $\Delta_{\mathrm{nd}}(C)=\infty$ means uncontrolled transport.
\end{definition}

\begin{remark}[Status and dependency boundary]
This definition is provisional and interface-level.
It does not select a finalized operator realization and does not assert that the chosen gauges are canonical.
Its current role is to provide a clean first analytic control target for \cref{lem:stage6-first-concrete-no-drift}, downstream of Stage-5 local uniqueness and upstream of \cref{thm:stage6-neutralizing-transport-uniformity}.
\end{remark}
